\section{Data Structures and Core Types}

To establish the tracking mechanism, several foundational data types and state variables had to be introduced.

These modifications are all contained within several of the project's main headers: \texttt{include/types.h}, \texttt{include/afl-fuzz.h}, \texttt{include/sharedmem.h}, and \texttt{include/config.h}.

\subsubsection{Configurable Slot Types}

To provide tunable fidelity, the width of the bitfield used to track distinct targets originating from a single indirect source --- i.e., the size of each Bloom filter --- is made configurable via the compiler preprocessor. This configuration is codified by the \texttt{INDIR\_SLOT\_SIZE} compile-time macro. It defaults to a size of 64.

The header \texttt{types.h} subsequently constructs the \texttt{indir\_slot\_t} alias utilising preprocessor conditionals based on the value of the aforementioned macro:
\begin{itemize}
    \item \texttt{INDIR\_SLOT\_SIZE == 32}: \texttt{indir\_slot\_t} resolves to \texttt{uint32\_t}.
    
    \item \texttt{INDIR\_SLOT\_SIZE == 64}: \texttt{indir\_slot\_t} resolves to \texttt{uint64\_t}.
    
    \item \texttt{INDIR\_SLOT\_SIZE == 128}: \texttt{indir\_slot\_t} resolves to \texttt{unsigned \_\_int128}.
\end{itemize}

By exclusively employing \texttt{indir\_slot\_t} across the fuzzer, the entire architecture scales to accommodate the designated slot width without requiring localised code modifications. This ensures type safety and prevents hardcoded integer sizes throughout the codebase.

Alongside the type alias, the header defines the \texttt{INDIR\_BIT\_SHIFT} constant, representing the base 2 logarithm of the slot size --- e.g., 6 for 64-bit. This is used inside the callback to construct a bitmask used for the modulo reduction of the hashed target address.

\subsubsection{Fuzzer State Encapsulation}

The fuzzer's context is entirely encapsulated within the \texttt{afl\_state\_t} structure. This structure was extended to manage the new coverage domain. Additions include: 
\begin{itemize}
    \item \texttt{indir\_map}: A pointer to the shared memory region populated by the target binary.
    
    \item \texttt{indir\_virgin\_bits}: A persistent array storing the cumulative history of all discovered indirect path transitions across the entire campaign. This is the reference used to evaluate novelty after executions.
    
    \item \textit{Top Entries Structures}: Additional arrays were introduced to track and retain the inputs that maximise specific bit transitions within the indirect map, mirroring the behaviour of the primary map's top-rating system.
\end{itemize}

\subsubsection{Shared Memory IPC Metadata}

The \texttt{sharedmem\_t} structure maintains the IPC primitives and metadata required to manage the shared memory segments. The structure was extended with: 
\begin{itemize}
    \item \textbf{\texttt{indir\_map\_size}}: An integer recording the currently allocated size of the secondary shared memory segment. This is updated whenever the fuzzer issues a reallocation based on forkserver feedback.
    
    \item \textbf{\texttt{indir\_shm\_id}}: The active System V IPC identifier associated with the allocated \texttt{indir\_map}.
\end{itemize}